\chapter{مفاهيم أساسية} 
\section[المعادلة التفاضلية]{المعادلة التفاضلية}


    هي علاقة بين المتغير المعتمد والمتغير المستقل (المتغيرات المستقلة) تدخل فيها المشتقات أو التفاضلات وتسمى المعادلة التفاضلية اعتيادية إذا كان المتغير المعتمد دالة في متغير مستقل واحد وبالتالي لا تحتوي إلا على مشتقات عادية. حيث الشكل العام لها
\[
F(x,y',y'',y''',\dots,y^{(n)})=0
\]
حيث $y^{(n)}=\mfrac{d^ny}{dx^n}$ المشتقة من الرتبة $n$ للمتغير $y$ بالنسبة إلى $x$.

\begin{example}
    ليكن $x$ هو المتغير المستقل و $y$ هو المتغير المعتمد. فالعلاقات التالية تمثل معادلات تفاضلية اعتيادية
    \begin{align}
        &\frac{dy}{dx}+y=3\sin x\label{diff_eq_ex1}\\
        &\frac{d^2y}{dx^2}+\frac{1}{x}\frac{dy}{dx}+y=0\label{diff_eq_ex2}
    \end{align}
\end{example}

\section[رتبة المعادلة التفاضلية]{رتبة المعادلة التفاضلية}
    إذا كانت المشتقة النونية $y^{(n)}$ هي أعلى مشتقة تظهر في المعادلة التفاضلية الاعتيادية قيل أن هذه المعادلة من الرتبة $n$.

\begin{example}
    المعادلة \eqref{diff_eq_ex1} هي معادلة اعتيادية من الرتبة الأولى. أما المعادلة \eqref{diff_eq_ex2} هي معادلة تفاضلية اعتيادية من الرتبة الثانية.
\end{example}

\section[درجة المعادلة التفاضلية]{درجة المعادلة التفاضلية}
    هي الاس المرفوع إليها أعلى مشتقة تظهر في المعادلة التفاضلية ، وقبل تحديد درجة المعادلة التفاضلية يجب وضعها في أبسط صورة قياسية صحيحة من حيث المشتقات.
\newpage
\begin{example}
    المعادلة
    \[
    \pbracket{\frac{d^2y}{dx^2}}^3+x\pbracket{\frac{dy}{dx}}+x^2y^3=0
    \]
    هذه معادلة تفاضلية من الرتبة الثانية و الدرجة الثالثة. أما المعادلة 
    \[
    \sbracket{1+\pbracket{\frac{dy}{dx}}^3}^{1/2}+3\frac{d^2 y}{dx^2}+xy=0
    \]
    قبل تحديد درجة المعادلة يجب وضعها على صورة خالية من الجذور. بإجراء عمليات بسيطة نلاحظ أن
    \[
    \sbracket{1+\pbracket{\frac{dy}{dx}}^3}^{1/2}=-\left[3\pbracket{\frac{d^2 y}{dx^2}}+xy\right]
    \]
    بتربيع الطرفين
    \[
    1+\pbracket{\frac{dy}{dx}}^3=\pbracket{3\frac{d^2y}{dx^2}+xy}^2
    \]
    بإستخدام مربع حدانية للطرف الايمن ، نحصل على
    \[
    1+ \left(\frac{dy}{dx}\right)^3 = 9 \left(\frac{d^2 y}{dx^2}\right)^2 + 6xy \frac{d^2 y}{dx^2} + x^2 y^2
    \]
    بإعادة ترتيب الحدود ، نحصل على
    \begin{gather*} 9\pbracket{\frac{d^2y}{dx^2}}^2+6xy\frac{d^2y}{dx^2}-\pbracket{\frac{dy}{dx}}^2+x^2y^2-1=0
    \end{gather*}
    وهذه معادلة تفاضلية من الرتبة الثانية ومن الدرجة الثانية
\end{example}

\section[المعادلة التفاضلية الخطية]{المعادلة التفاضلية الخطية}
    هي المعادلة الخطية في المتغير المعتمد و مشتقاته جميعاً. الصورة العامة للمعادلة التفاضلية الخطية من الرتبة $n$ 
    \[
    p_n(x)y^{(n)}+p_{n-1}(x)y^{(n-1)}+\dots+p_1(x)y'+p_0(x)y=Q(x)
    \]
    أو
    \[
    \sum_{i=0}^{n}p_i(x)y^{(i)}=Q(x)
    \]
    حيث المتغير المعتمد $y$ و جميع مشتقاته مرفوعة إلى الاس واحد ولا توجد حواصل ضرب مشتركة في ما بينها ، و الدوال $Q(x)$ و $p_i(x)$ هي دوال للمتغير المستقل $x$ خطية أم غير خطية لا تؤثر على خطية المعادلة التفاضلية.

\begin{note}
    من تعريف المعادلة التفاضلية الخطية يمكننا القول أي معادلة تفاضلية فيها المتغير المعتمد $y$ أو أحد مشتقاته مرفوعة إلى اس غير الواحد أو وجدنا حاصل ضرب في ما بينها. سوف نطلق عليها معادلة تفاضلية غير خطية.
\end{note}

\begin{example}
    لدينا المعادلة
    \[
    x^2y^{''}+xy'+y=\cos x
    \]
    معادلة تفاضلية اعتيادية خطية. أما المعادلة
    \[
    yy^{'}+y^{''}=e^{3x}
    \]
    هذه المعادلة التفاضلية غير خطية لوجود حاصل الضرب $yy^{'}$. وأيضا لدنيا
\[
y^{'}+x\sqrt[4]{y}=\cot x
\]
وهذه كذلك ليست خطية فيها المتغير المعتمد $y$ مرفوع إلى الاس $1/4$ (غير الواحد).
\end{example}

\section[المعادلة التفاضلية الخطية المتجانسة]{المعادلة التفاضلية الخطية المتجانسة}
اذا انعدمت الدالة $Q(x)$ (اي ان $Q(x) = 0$) من المعادلة التفاضلية
\[
p_n(x)y^{(n)}+p_{n-1}(x)y^{(n-1)}+\dots+p_1(x)y'+p_0(x)y=Q(x)
\]
أو
\[
\sum_{i=0}^{n}p_i(x)y^{(i)}=Q(x)
\]
لجميع قيم $x$ فأنها تصبح معادلة خطية متجانسة ، عدا ذلك تكون المعادلة غير متجانسة

\begin{example}
	المعادلة
	\[
	x^2 y'' + xy' + (x^2 - 1)y = 0
	\]
	هي معادلة تفاضلية خطية متجانسة من الرتبة الثانية. بينما المعادلة
	\[
	xy' + (\sin x)y = x^2(\sin x + 2)
	\] 
	هي معادلة تفاضلية خطية غير متجانسة من الرتبة الاولى.
\end{example}

\section[حل المعادلة التفاضلية]{حل المعادلة التفاضلية}
تسمى الدالة $y=y(x)$ حلاً للمعادلة التفاضلية $F(x,y,y',y'',\dots,y^{(n)})$ إذا كانت
\begin{enumerate}
    \item قابلة للاشتقاق $n$ من المرات
    \item تحقق المعادلة التفاضلية أي أن $F(x,y(x),y'(x),y''(x),\dots,y^{(n)}(x))$
\end{enumerate}
\begin{example}
    أثبت أن $y(x)=c \sin x$ حلاً للمعادلة التفاضلية $y''+y=0$ حيث $c$ ثابت.
\end{example}
\begin{solution}
    نشتق العلاقة ونقوم بتعويضها في المعادلة التفاضلية
\[
y=c\sin x\Rightarrow y'=c\cos x\Rightarrow y''=-c\sin x
\]
الآن نعوض
\[
y''+y=-c\sin x+c\sin x=0
\]
\end{solution}
\section[الحل العام و الحل الخاص]{الحل العام و الحل الخاص}
الحل العام للمعادلة التفاضلية من الرتبة $n$ هو حل يحتوي على $n$ من الثوابت الاختيارية وبالطبع يحقق المعادلة التفاضلية.\\
أما الحل الخاص هو أي حل يحقق المعادلة التفاضلية لا يشتمل على أي ثوابت اختيارية وقد نحصل عليه احيانا بالتعويض عن الثوابت الاختيارية في الحل العام بقيم محددة.

\begin{example}
    الحل العام للمعادلة التفاضلية $y'''-5y''+6y'=0$ يمكن ايجادمن خلال المعادلة المميزة للمعادلة التفاضلية الناتجة من تحويل شكل المعادلة الى صيغة المؤثر
    \[
    (D^3 - 5D^2 + 6D) y = 0
    \]
    اذن المعادلة المميزة تكون
    \[
    m^3 - 5m^2 + 6m = 0
    \]
    نلاحظ يمكننا ان نحلل الطرف الايسر بالعامل المشترك
    \[
    m(m^2 - 5m + 6) = 0
    \]
    الآن بالتجربة
    \[
    m(m-2)(m-3) = 0
    \]
    اذن الحلول للمعادلة المميزة
    \[
    m_1 = 0 , m_2 = 2, m_3 = 3
    \]
    بالتالي الحل العام يكون
    \[
    y = C_1 e^{m_1x} + C_2 e^{m_2 x} + C_3 e^{m_3 x}
    \]
   حيث $C_1, C_2, C_3 $ ثوابت اختيارية ، اذن الحل العام النهائي
    \[
    y(x) = C_1 + C_2 e^{2x} + C_3 e^{3x}
    \]
     ويمكننا ان نختار قيم محددة للثوابت الاختيارية ويتحول الحل العام الى حل خاص ، مثلاً نختار \linebreak
     $C_1 = 1, C_2, -4, C_3 = 3$ ، ليصبح الحل
    \[
    y(x) = 1 -4e^{2x} + 3e^{3x}
    \]
\end{example}

\begin{example}\label{ex:popexample}
	اوجد حل المعادلة التالية
	\[
	p' = 1.1 p
	\]
	مع الشرط الابتدائي $p(0)=2$
\end{example}
\begin{solution}
	نفصل المتغيرات فنكتب
	\[
	\frac{dp}{p} = 1.1\, dt
	\]
	ثم نكامل الطرفين فنحصل على
	\[
	\ln |p| = 1.1 t + C
	\]
	وبأخذ الأس للطرفين ينتج
	\[
	p = C e^{1.1 t}
	\]
	وباستخدام الشرط الابتدائي \(p(0)=2\) نجد أن \(C=2\)، وبالتالي يكون الحل
	\[
	p(t) = 2 e^{1.1 t}.
	\]
\end{solution}

\begin{example}
	اوجد حل المعادلة التالية
	\[
	y' +2y = \cos x
	\]
	مع الشرط الابتدائي $y(0)=1 $.
\end{example}
\begin{solution}
	نلاحظ أن المعادلة خطية من الدرجة الأولى، فنستخدم عامل التكامل \(e^{\int 2\,dx}=e^{2x}\)، وبضرب طرفي المعادلة به نحصل على
	\[
	e^{2x}y' + 2e^{2x}y = e^{2x}\cos x
	\]
	والطرف الأيسر يمثل مشتقة حاصل الضرب:
	\[
	\frac{d}{dx}\big(e^{2x}y\big) = e^{2x}\cos x
	\]
	بالتكامل نحصل على
	\[
	e^{2x}y = \int e^{2x}\cos x\,dx
	\]
	ومن التكامل المعروف:
	\[
	\int e^{2x}\cos x\,dx = \frac{e^{2x}(2\cos x + \sin x)}{5} + C
	\]
	إذن
	\[
	e^{2x}y = \frac{e^{2x}(2\cos x + \sin x)}{5} + C
	\]
	وبالقسمة على \(e^{2x}\):
	\[
	y = \frac{2\cos x + \sin x}{5} + Ce^{-2x}
	\]
	وباستخدام الشرط الابتدائي \(y(0)=1\):
	\[
	1 = \frac{2}{5} + C
	\Rightarrow C = \frac{3}{5}
	\]
	وعليه يكون الحل:
	\[
	y(x) = \frac{2\cos x + \sin x}{5} + \frac{3}{5}e^{-2x}.
	\]
\end{solution}

\section{طريقة رانج كتا من الرتبة الرابعة \en{RK4}}

تُستخدم طريقة رانج كتا من الرتبة الرابعة لحل المعادلات التفاضلية العادية من الشكل:
\[
\frac{dy}{dx} = f(x, y), \quad y(x_0) = y_0
\]
نفترض أننا نريد إيجاد الحل على الفترة \( [x_0, x_n] \)، فنقوم بتقسيمها إلى \( n \) أجزاء متساوية:
\[
h = \frac{x_n - x_0}{n}
\]
ونحصل على نقاط التقسيم:
\[
x_1 = x_0 + h,\quad x_2 = x_0 + 2h,\quad \dots,\quad x_n = x_0 + nh
\]
وتهدف طريقة رانج كتا لإيجاد قيم $y$ عند قيم $x$، وهي كالتالي:
\begin{itemize}
	\item \( y_0 \): القيمة الابتدائية المعطاة عند \( x_0 \).
	\item \( y_1 \): تقريب لقيمة الحل عند النقطة \( x_1 \).
	\item \( y_2 \): تقريب لقيمة الحل عند \( x_2 \).
\end{itemize}
وهكذا ، 
أي أن \( y_k \) يمثل تقريبًا للقيمة الحقيقية للدالة \( y(x_k) \)، ويتم حساب هذه القيم بشكل تتابعي باستخدام الطريقة.
لإيجاد \( y_{n+1} \) من \( y_n \) باستخدام خطوة \( h \)، نحسب:
\[
k_1 = f(x_n, y_n)
\]
\[
k_2 = f\left(x_n + \frac{h}{2}, \, y_n + \frac{h}{2} k_1 \right)
\]
\[
k_3 = f\left(x_n + \frac{h}{2}, \, y_n + \frac{h}{2} k_2 \right)
\]
\[
k_4 = f(x_n + h, \, y_n + h k_3)
\]
ثم نحسب القيمة الجديدة:
\[
y_{n+1} = y_n + \frac{h}{6} \left( k_1 + 2k_2 + 2k_3 + k_4 \right)
\]